\documentclass[10pt,a4paper]{amsart}
\usepackage[margin=19mm]{geometry}
\usepackage{amsmath,amssymb,mathtools}
\usepackage[scr=rsfs,cal=euler]{mathalfa}
\usepackage{xfrac}
\usepackage[unicode,hidelinks,bookmarks=false]{hyperref}
\newcommand{\ie}{i.e.\ }
\newcommand{\cf}{cf.\ }
\newcommand{\resp}{respectively\ }
\newcommand{\Subsection}{\subsection}
\providecommand{\nobreakdash}{\nobreak\hbox{-}}
\setlength{\emergencystretch}{2em}
\multlinegap0pt
\usepackage{bm}
\usepackage{bbm}
\usepackage{dsfont}
\usepackage{xspace}
\usepackage{cancel}
\usepackage{mathtools}
\usepackage{cleveref}
\usepackage[shortlabels]{enumitem}
\usepackage{amsthm}
\makeatletter
\@ifundefined{thm}{\newtheorem{thm}{Thm}}{}
\@ifundefined{claim}{\newtheorem{claim}[thm]{Claim}}{}
\@ifundefined{claimproof}{\newtheorem{claimproof}{Claimproof}}{}
\@ifundefined{defn}{\newtheorem{defn}[thm]{Definition}}{}
\@ifundefined{prop}{\newtheorem{prop}[thm]{Proposition}}{}
\@ifundefined{rem}{\newtheorem{rem}[thm]{Remark}}{}
\makeatother
\newcommand{\C}[1]{\mathcal{C}(#1)}
\newcommand{\Hor}[1]{\operatorname{Hor}\left(#1\right)}
\newcommand{\N}{\mathbb{N}}
\newcommand{\W}{\mathcal{W}}
\newcommand{\depth}[1]{\operatorname{Depth}\left(#1\right)}
\newcommand{\diam}{\mathrm{diam}}
\newcommand{\dist}{\mathrm{d}}
\newcommand{\down}[1]{{#1}^{\downarrow}}
\newcommand{\frakS}{\mathfrak S}
\newcommand{\link}[2]{\operatorname{Link}_{#1}\left(#2\right)}
\newcommand{\ov}[1]{\overline{#1}}
\newcommand{\squid}[1]{squid}
\renewcommand{\hat}{\widehat}
\title[Cusped spaces for hierarchically hyperbolic groups, and appl]{Cusped spaces for hierarchically hyperbolic groups, and applications to Dehn filling quotients}
\author{Giorgio Mangioni, Alessandro Sisto}
\date{Source: arXiv:2602.23275v2 (CC BY 4.0)}
\begin{document}
\maketitle

\noindent\emph{Source and licence.} Cusped spaces for hierarchically hyperbolic groups, and applications to Dehn filling quotients. Version arXiv:2602.23275v2 (CC BY 4.0), distributed under the Creative Commons Attribution 4.0 International licence, as recorded in the arXiv licence metadata for this exact version and shown on the abstract page https://arxiv.org/abs/2602.23275v2. Full rights notice: licenses/arxiv-2602.23275-CC-BY-4.0.txt.\par\smallskip

\noindent\textit{Editorial scope.} Counted unit: the Prop whose source label is \texttt{prop\_coarseemb} in the article (source0.tex lines 574--578, source label \texttt{prop\_coarseemb}), delivered together with the article's own complete original proof of it (source0.tex lines 581--595, one \texttt{proof} environment of 15 lines). No mathematics is written, repaired or re-derived: statement and proof are the article's own text line for line. Required context retained uncounted: def:horoball (lines 218--225); CHHS\_def\_edge\_in\_link (lines 284--290); rem:depth\_diff (lines 550--552); item\_coarseemb (lines 560--565). Every definition, display and label of those lines is retained unchanged. Omitted from this package and not redistributed: 1--217, 226--283, 291--549, 553--559, 566--573, 579--580, 596--1364. Nothing inside a retained excerpt is omitted or altered: every retained body is delivered exactly as the article's lines are written, so no omission receipt of any kind accompanies this package. Citations: the retained excerpts carry no citation command at all, so no bibliography entry of the article is transcribed and no reference is redistributed. Reference adaptation: none was needed and none was made; every reference inside the delivered unit resolves inside it, so no unresolved reference or citation is produced. Printed slips kept literal: no printed word, symbol, hypothesis or number of the counted statement or of its proof was altered. Layout and class: the article's own class is \texttt{amsart} with packages including babel, inputenc, geometry, hyperref, cleveref, amssymb, amsfonts, amsmath, amsthm, amscd, amsxtra, latexsym, mathabx, mathtools; the wrapper is the lane's pinned \texttt{amsart} editorial wrapper at 10pt on A4, which restates the theorem-like environments the retained text needs and adds the article's own macro definitions that the retained text needs (\texttt{\textbackslash C}, \texttt{\textbackslash Hor}, \texttt{\textbackslash N}, \texttt{\textbackslash W}, \texttt{\textbackslash depth}, \texttt{\textbackslash diam}, \texttt{\textbackslash dist}, \texttt{\textbackslash down}, \texttt{\textbackslash frakS}, \texttt{\textbackslash hat}, \texttt{\textbackslash link}, \texttt{\textbackslash ov}, \texttt{\textbackslash squid}), with extra wrapper packages (bm, bbm, dsfont, xspace, cancel, mathtools, cleveref, enumitem). The delivered numbering is the unit's own. Assets: no retained span contains a figure environment, a graphics-inclusion command, a PDF-page inclusion, a table or a TikZ picture; no image file is copied into this package, the package declares no asset dependency and no image file is redistributed with it. Licence: version arXiv:2602.23275v2 of this article is distributed under the Creative Commons Attribution 4.0 International licence, as recorded in the arXiv licence metadata for this exact version and shown on the abstract page https://arxiv.org/abs/2602.23275v2. A full rights notice is delivered at licenses/arxiv-2602.23275-CC-BY-4.0.txt. Characters: every retained excerpt is pure ASCII, so the delivered engine, XeLaTeX with its default Latin Modern text font, reports no missing character.
\begin{defn}\label{def:horoball}
Let $\Gamma$ be a simplicial graph. The \emph{combinatorial horoball} $\Hor{\Gamma}$ with base $\Gamma$ is the simplicial graph whose vertex set is $\Gamma^{(0)}\times \N$, and with the following two types of edges:
\begin{itemize}
\item For every $p\in \Gamma^{(0)}$ and $n\in \N$, $(p,n)$ is adjacent to $(p, n+1)$;
\item For every $p,q\in \Gamma^{(0)}$ and $n\in \N$, if $\dist_{\Gamma}(p,q)\le 2^n$ then $(p,n)$ is adjacent to $(q,n)$.
\end{itemize}
We often denote $(p,n)\in \Gamma^{(0)}\times \N$ simply by $p^n$. 
\end{defn}

\begin{enumerate}[label=(\arabic*)]
\item \label{CHHS_def_complexity} Any chain $\link{X}{\Delta_1}\subsetneq\dots\subsetneq\link{X}{\Delta_k}$ has length at most $n$, which is called the \emph{complexity} of $(X,\W)$.
\item \label{CHHS_def_hyperbolicity} For each $[\Sigma]\in \frakS$,  $\C{\Sigma}$ is $\delta$--hyperbolic.
\item \label{CHHS_def_qi_emb} For each $[\Sigma]\in \frakS$, $\C{\Sigma}$ is $\delta$--quasi-isometrically embedded in $Y_\Sigma$.
\item \label{CHHS_def_lk_cap_lk} For every $[\Delta],[\Sigma]\in \frakS$ for which there exists $[\Gamma]\in\frakS$ such that $\link{X}{\Gamma}\subseteq \link{X}{\Delta}\cap \link{X}{\Sigma}$ and $\diam(\C{\Gamma})\geq \delta$, there exists a non-maximal simplex $\Pi$ containing $\Sigma$ such that $\link{X}{\Pi}\subseteq \link{X}{\Delta}$, and all $\Gamma$ as above satisfy $\link{X}{\Gamma}\subseteq \link{X}{\Pi}$.
\item \label{CHHS_def_edge_in_link} For each $[\Delta]\in \frakS$ and $v\neq w\in \link{X}{\Delta}$, if $v,w$ are adjacent in $\C{\Delta}$ but not in $\link{X}{\Delta}$, then there exist $\mathcal W$-adjacent maximal simplices $\Pi_v,\Pi_w$ such that $\Delta\star \{v\}\subseteq \Pi_v$ and $\Delta\star \{w\}\subseteq \Pi_w$.
\end{enumerate}

\begin{rem}\label{rem:depth_diff}
   If $\Theta,\Xi\in \hat \W^{(0)}$ are $\hat \W$-adjacent, then $|\depth{\Theta}-\depth{\Xi}|\le 1$ by construction.
\end{rem}

\begin{enumerate}[label=(\Alph*)]
\item\label{item_coarseemb} The inclusion $X\hookrightarrow \hat X$ induces an injective coarse embedding $i\colon\W\hookrightarrow \hat \W$. Moreover, the coarse embedding functions do not depend on any of the data.
\item\label{item_Teich_is_combhhs} There exists $\delta'$, only depending on $\delta$, such that $(\hat X,\hat \W, n,\delta')$ is a squid HHS.
\item\label{item_G-action_on_teich} Any action of a group $G$ on $(X,\W)$ extends to a $G$-action on $(\hat X,\hat \W)$.
\item\label{item_W_is_relHHS} $\W$ is a hierarchically hyperbolic space relative to $\{\C{\Delta_v}\}_{v\in \ov X^{(0)}}$.
\end{enumerate}

\begin{prop}\label{prop_coarseemb}
\ref{item_coarseemb} Let $(X,\W)$ be a relative \squid{} HHS. The inclusion $X\hookrightarrow \hat X$ induces an injective coarse embedding $\W\hookrightarrow \hat \W$. More precisely, for every two maximal simplices $\Theta,\Xi\in \W^{(0)}$, 
$$m(\dist_{\W}(\Theta,\Xi))\le \dist_{\hat \W}(\Theta,\Xi)\le \dist_{\W}(\Theta,\Xi),$$
where $m$ is the inverse\footnote{This function is known as the Lambert $W$ function, and diverges at infinity. See e.g. \href{https://mathworld.wolfram.com/LambertW-Function.html}{https://mathworld.wolfram.com/LambertW-Function.html} for more details.} of the function $n\to n2^n$.
\end{prop}

\begin{proof}
    Let $\Theta$ be a maximal simplex of $X$, which we can also see as a maximal simplex of $\hat X$ with $\Theta^+=\emptyset$ and $\down{\Theta}=\Theta$. Thus there is an injection $\W\hookrightarrow \hat \W$, which maps $\W$-edges to $\hat \W$-edges of $\W$-type. Since $\W$ is an induced subgraph of $\hat \W$, the inclusion $\W\hookrightarrow\hat \W$ is $1$-Lipschitz. 
    
    We are left to prove that, if $\Theta,\Xi\in \W^{(0)}$ and $D=\dist_{\hat \W}(\Theta,\Xi)$, then $\dist_{\W}(\Theta, \Xi)\le D2^D.$ To this extent, let $\{\Sigma_0=\Theta, \Sigma_1,\ldots, \Sigma_D=\Xi\}$ be a $\hat \W$-geodesic, and consider the path $\{\down{\Sigma_0}=\Theta, \down{\Sigma_1},\ldots, \down{\Sigma_D}=\Xi\}$ in $\W$.
    
    \begin{claim}\label{claim:coarseemb}
        For every $i=0,\ldots, D-1$, $\dist_{W}(\down{\Sigma_i}, \down{\Sigma_{i+1}})\le 2^D$.
    \end{claim}

    \begin{claimproof}[Proof of \Cref{claim:coarseemb}]
        If $\Sigma_i$ and $\Sigma_{i+1}$ span an edge of $\W$-type then $\down{\Sigma_i}$ and $\down{\Sigma_{i+1}}$ are $\W$-adjacent, and therefore  $\dist_{W}(\down{\Sigma_i}, \down{\Sigma_{i+1}})\le 1$. Thus suppose that $\Sigma_i$ and $\Sigma_{i+1}$ are joined by an edge of cusp-type, so that there exist $v\in \ov X^{(0)}$ and $\Hor{\C{\down{\Delta_v}}}$-adjacent points $p^n,q^m\in L_v\times \N$ such that $\Sigma_i-\Sigma_{i+1}=\{p^n\}$ and $\Sigma_{i+1}-\Sigma_{i}=\{q^m\}$. By inspection of \Cref{def:horoball} there are two sub-cases to analyse. If $q^m=p^{n+1}$ then $\down{\Sigma_{i}}= \down{\Sigma_{i+1}}$, and we are done. Otherwise we must have that $m=n$ and $\dist_{\C{\Delta_v}}(p,q)\le 2^n$, so there is a $\C{\Delta_v}$-geodesic with vertices $p_0=p,p_1,\ldots,p_k=q$ for some $k\le 2^n$. Since $L_v\times \N=\link{X}{\Sigma_i\cap\Sigma_{i+1}}$, \Cref{CHHS_def_edge_in_link} for $(X,\W)$ then implies that $\Sigma_i$ and $\Sigma_{i+1}$ are at distance at most $k\le 2^n$ in $\W$. But now $n\le \depth{\Sigma_i}$, which by \Cref{rem:depth_diff} is at most $D+\depth{\Sigma_0}=D$, as required.
    \end{claimproof}
    \noindent \Cref{prop_coarseemb} now follows from the Claim, since 
    \begin{align*}\dist_{W}(\Theta, \Xi)\le \sum_{i+0}^{D-1}\dist_{W}(\Sigma_i, \Sigma_{i+1})\le D2^D. & \qedhere \end{align*}
\end{proof}

\end{document}
